\section{Verification record}\label{si:verification}

\subsection{Claim layers and evidence}

\begin{table}[!htbp]
\centering\small
\caption{What supports the main classes of claims, and what it does not establish.}
\begin{tabularx}{\textwidth}{@{}>{\raggedright\arraybackslash}p{0.24\textwidth}XX@{}}
\toprule
Claim layer & Evidence & Limitation \\
\midrule
Continuum identities for the witness & Written proofs in the main appendices and SI~S2--S4 & These are paper proofs; they are not established by finite grid checks \\
Finite witness classifications & Exact-integer grid artifact in Figshare v3 (25 checks on a $31^3$ state grid at scale 40) & Verifies only the enumerated rational grid and specified predicates \\
Benchmark arithmetic & 24 exact-rational checks in the SafeTransition v1 artifact & Checks the declared numeric illustration; it does not supply an empirical income model or calibrate a cod stock \\
Selected CES point signs & Four rational logarithm enclosures and a 361-point harmonic identity check reported in the supplementary audit & The scripts/full bounds were not located in the inspected DOI packages; finite samples do not prove a continuum landscape \\
CES continuum statements & Diagonal and pointwise formulas proved in SI~S2--S3 & No global fixed-sum stationary-point, equal-minimum, or slice-closure theorem is claimed in v77 \\
Lean formalization & Pinned source contains selected algebraic witness components and CES rung examples & It does not formalize the full real-exponent nesting, continuity/IVT transfer, or relative-interior argument; no v77 Lean build was run \\
\bottomrule
\end{tabularx}
\end{table}

\subsection{Finite verifiers and their scope}

The Figshare v3 exact-integer verifier enumerates $29{,}791$ states $(x,s_1,s_2)$ on a scale-40 grid and records 25 checks of witness classifications, region identities, licensing thresholds, endpoint comparisons, and hold-prefix persistence. This is exact finite arithmetic on the listed grid, not a proof over the continuum. The SafeTransition v1 benchmark script reports 24 rational checks for the model quantities used in the earlier benchmark package and the abstract action tubes in main-text Table~\ref{main-tab:actions}; v77 changes the unit interpretation and states explicitly that the income coordinate is stipulated, not empirically derived.

The earlier report also listed separate CES, off-diagonal, continuum, and common-shock check counts (57, 50, 44, and 77, respectively). Their scripts and run logs were not present in the two inspected Figshare deposits. A supplementary audit script was rerun with its label-scan input redirected to the current modular SI; that adapted run reproduced four rationally enclosed geometric point verdicts, 361 exact rational harmonic checks, and a numerical tangency search. The geometric verdicts are finite point certificates; the tangency values are high-precision numerical estimates, not certified roots or global extrema. The former 44-check list is not a continuum certificate and is not used to support any global fixed-sum claim in v77. The audit's number of states or checks must not be read as an independent proof of a continuum theorem.

\subsection{Lean coverage and limits}

The cited Lean source is pinned at commit \texttt{6f626ee4\allowbreak 8b2a3a59\allowbreak 282b5df5\allowbreak adda76c8\allowbreak d5993fdc}: \href{https://github.com/MIKEAA2020/general-sustainability/blob/6f626ee48b2a3a59282b5df5adda76c8d5993fdc/lean/Formalizations/P1_AssessmentSeparation.lean}{pinned source file}. Its header explicitly excludes the nesting of real-exponent power means, transfer to non-ladder elasticities, the continuity/IVT step, and relative-interior topology (replaced there by explicit witnesses). The module includes selected algebraic witness identities and rung-specific CES statements, but this does not amount to formalization of the paper's entire pointwise and all-positive-elasticity continuum theorem. The count of declarations is not a coverage metric. No Lean build, dependency-axiom audit, or real-number instantiation check was run for v77; those statuses are not asserted here.

\subsection{Reproduction and version status}

The v77 local package includes assembled and modular TeX, figures, the modified Figure~\ref{main-fig:triangle} source script, a source script for the benchmark-axis unit-label edit, build instructions, hashes, and local QA records. It has no DOI and has not been externally deposited. The existing Figshare and SafeTransition records are version-mismatched; no single inspected deposit reproduces every reported calculation from v77. The main-text Data Availability statement lists those boundaries. The DFO historical anchor was independently traced to Table A2, PDF p.~16, of CSAS SAR 2016/026: $6192.06$ kt over 1983--1989 gives a mean of $884.58$ kt, rounded to $884.6$ kt; the 2015 value is approximately $298.65$ kt. This is a historical assessment-vintage check, not a current threshold update.

\section{Positioning notes, per literature}\label{si:positioning}

\paragraph{Commensurability} The foundation statement that values may be only weakly comparable -- and that this incommensurability is constitutive of ecological economics -- is due to Martinez-Alier, Munda and O'Neill (1998). The main article's operators give one formal model of the distinction that paper draws between strong and weak commensurability.

\paragraph{Maximin and indicator formulations} The characterization of sustainability indicators and the role of maximin and MSY-type reasoning in operationalizing strong-sustainability criteria are developed in Martinet (2011), Cairns and Martinet (2014), and Doyen and Gajardo (2020). This literature includes formal viability and Pareto formulations; the typed operator is one narrow transition-certification idealization and is not presented as a comprehensive or uniquely strongest formalization.

\paragraph{Compensability in multi-criteria decision analysis} Compensability analysis is established in MCDA, including mappings between compensatory aggregation, weak sustainability, and outranking methods (Cinelli, Coles and Kirwan, 2014; Sch\"ar, Pohl and Geldermann, 2025). Weighted-sum limitations on nonconvex Pareto fronts are also well known (Das and Dennis, 1997). The present finite-menu theorem shares convex-hull geometry with these results; it asks a different question about state acceptance and plan selection, and the same quantifier gap can occur for static alternatives.

\paragraph{Acceptance semantics and decision aiding} The protocols of main-text Table~\ref{main-tab:protocols} are decision-aiding constructs: each specifies the rule by which a composite reading sorts a transition into the accepted or rejected category -- the sorting problematic of decision aiding (Roy, 1996; Vincke, 1992). Under the weight-adaptive rule a transition is accepted when every weighting licenses some plan, possibly a different plan for each weighting; the licensing thresholds are then the verdict's weight-sensitivity intervals, analogous to sensitivity analyses applied to composite-indicator rankings (Nardo et al., 2008). The preference-elicitation background is multi-attribute value theory's substitution rates (Keeney and Raiffa, 1976) and the incommensurable-criteria treatment of sustainable development (Munda, 2005).

\paragraph{Viability and robust control} The backward recursion is a typed instance of robust-predecessor, reach--avoid, capture-basin and hybrid-reachability constructions (Aubin, 1991; Aubin, Bayen and Saint-Pierre, 2011; Saint-Pierre, 1994; Lygeros, Tomlin and Sastry, 1999; Aubin and Catt\'e, 2002), with natural-resource management as an established application (De Lara and Doyen, 2008; Doyen et al., 2012). Fractional physical blending and complete-excursion concatenation compare distinct control classes; they are not a mixed/behavioral strategy equivalence. The here-and-now/wait-and-see analogy from adjustable robust optimization concerns the order of plan and weight quantifiers only; the weight is not modeled as a realized disturbance (Ben-Tal et al., 2004).

\paragraph{Planetary boundaries and the safe-and-just space} The typed operator is a formal idealisation of the separate-boundary reading underlying the Planetary Boundaries framework (Rockstr\"om et al., 2009; Steffen et al., 2015), and the same separately-checked-floor logic underlies the safe-and-just-space literature (Raworth, 2012; O'Neill et al., 2018; Fanning et al., 2022). Boundary processes interact and amplify human impacts (Lade et al., 2020), which is why the coupled disturbance regime of S4 is not a remote possibility. The main theorem provides a formal model of \emph{one reason} a noncompensatory stance may be adopted; it does not prove that planetary boundaries ought to be noncompensatory, which is a normative and scientific-design question.

\section{Finite-generation consequence under compact confinement}\label{si:finitegeneration}

This section retains one concise, fully specified consequence of the earlier intergenerational discussion. Let $K$ be a compact state set, and let $V^{(k)}\subseteq K$ be closed, nested decreasing generation constraints with $\bigcap_{k\ge1}V^{(k)}=\varnothing$.

\begin{proposition}[finite-generation obstruction]\label{prop:finitegeneration}
Under these assumptions there is a finite generation $N$ for which $V^{(N)}=\varnothing$. Consequently, no trajectory confined to $K$ can satisfy the requirements through every generation.
\end{proposition}
\begin{proof}
If every $V^{(k)}$ were nonempty, the nested family of closed subsets of compact $K$ would have the finite-intersection property. Compactness would imply $\bigcap_{k\ge1}V^{(k)}\ne\varnothing$, contrary to the hypothesis. Hence some $V^{(N)}$ is empty. Since the sets are decreasing, every later generation set is also empty.
\end{proof}

The common compact confinement is essential to this finite-generation conclusion. Without it, take constant dynamics on $\mathbb R$ and $V^{(k)}=[k,\infty)$. Every finite prefix through generation $N$ is feasible from $x=N$, but no state satisfies all generations. The initial state used for each finite prefix escapes every common compact set.

A separate monotonicity point is not conjectural: for fixed dynamics and a fixed policy class, if safe sets satisfy $K_1\subseteq K_2$, then the viability kernel obeys $\operatorname{Viab}(K_1)\subseteq\operatorname{Viab}(K_2)$. The earlier conjecture that expanding a safe boundary can contract that kernel by admitting a new attractor is therefore withdrawn under those fixed-model assumptions. Any different phenomenon would have to change the dynamics, objective, or policy class as well, and is outside v77.

\section*{References cited in the Supplementary Information}
\begin{itemize}[leftmargin=1.5em,itemsep=1pt]
\item Abaee, A. (2026a). An obstruction calculus for viability under incomplete observation. Zenodo. \url{https://doi.org/10.5281/zenodo.22552616}.
\item Abaee, A. (2026c). Does a surplus-production ladder improve forecasts of Northern cod? A scored test on NAFO 2J3KL. Zenodo. \url{https://doi.org/10.5281/zenodo.22553609}.
\item Abaee, A. (2026d). Verification code and figure pipeline for ``Aggregate indices and transition safety''. Figshare, version 3 (the attached manuscript is Paper1 v42, not this local Paper06 v77 edition). \url{https://doi.org/10.6084/m9.figshare.33764023.v3}.
\item Abaee, A. (2026e). Delay-induced regime change in harvested stocks: the mobilising and protective channels of institutional feedback, and the review interval as control. Zenodo. \url{https://doi.org/10.5281/zenodo.22554217}.
\item Abaee, A. (2026f). \emph{SafeTransition: exact rational certification of transition safety for sustainability assessment}. Software release 1.3.0, Figshare, version 1. \url{https://doi.org/10.6084/m9.figshare.33971389.v1}.
\item Aubin, J.-P. (1991). \emph{Viability Theory}. Birkh\"auser, Boston. Aubin, J.-P., Bayen, A.~M., and Saint-Pierre, P. (2011). \emph{Viability Theory: New Directions}, 2nd ed. Birkh\"auser. Aubin, J.-P., and Catt\'e, F. (2002). \emph{Set-Valued Analysis}, 10(4), 379--416.
\item Ben-Tal, A., Goryashko, A., Guslitzer, E., and Nemirovski, A. (2004). \emph{Mathematical Programming}, 99(2), 351--376.
\item Cairns, R.~D., and Martinet, V. (2014). \emph{European Economic Review}, 69, 4--17.
\item Cardaliaguet, P. (1996). \emph{SIAM Journal on Control and Optimization}, 34(4), 1441--1460.
\item Cinelli, M., Coles, S.~R., and Kirwan, K. (2014). \emph{Ecological Indicators}, 46, 138--148.
\item Das, I., and Dennis, J.~E. (1997). \emph{Structural Optimization}, 14, 63--69.
\item Dearing, J.~A., et al. (2014). Safe and just operating spaces for regional social-ecological systems. \emph{Global Environmental Change}, 28, 227--238.
\item De Lara, M., and Doyen, L. (2008). \emph{Sustainable Management of Natural Resources}. Springer. Doyen, L., et al. (2012). \emph{Ecological Economics}, 75, 32--42. Doyen, L., and Gajardo, P. (2020). \emph{Natural Resource Modeling}, 33(3), e12250.
\item DFO (2016). \emph{Stock Assessment of Northern Cod (NAFO Divs.\ 2J3KL) in 2016}. CSAS Science Advisory Report 2016/026.
\item Fanning, A.~L., O'Neill, D.~W., Hickel, J., and Roux, N. (2022). \emph{Nature Sustainability}, 5(1), 26--36.
\item Keeney, R.~L., and Raiffa, H. (1976). \emph{Decisions with Multiple Objectives}. Wiley.
\item Lade, S.~J., et al. (2020). \emph{Nature Sustainability}, 3(2), 119--128.
\item Lygeros, J., Tomlin, C., and Sastry, S. (1999). \emph{Automatica}, 35(3), 349--370.
\item Martinet, V. (2011). \emph{Journal of Environmental Economics and Management}, 61(2), 183--197.
\item Martinez-Alier, J., Munda, G., and O'Neill, J. (1998). \emph{Ecological Economics}, 26(3), 277--286.
\item Munda, G. (2005). In \emph{Multiple Criteria Decision Analysis: State of the Art Surveys}. Springer, 953--986.
\item Nardo, M., et al. (2008). \emph{Handbook on Constructing Composite Indicators}. OECD.
\item O'Neill, D.~W., Fanning, A.~L., Lamb, W.~F., and Steinberger, J.~K. (2018). \emph{Nature Sustainability}, 1(2), 88--95.
\item Raworth, K. (2012). \emph{Oxfam Policy and Practice: Climate Change and Resilience}, 8(1), 1--26.
\item Rockstr\"om, J., et al. (2009). \emph{Nature}, 461, 472--475. Steffen, W., et al. (2015). Planetary boundaries: Guiding human development on a changing planet. \emph{Science}, 347(6223), 1259855.
\item Roy, B. (1996). \emph{Multicriteria Methodology for Decision Aiding}. Kluwer. Vincke, P. (1992). \emph{Multicriteria Decision-Aid}. Wiley.
\item Saint-Pierre, P. (1994). \emph{Applied Mathematics and Optimization}, 29, 187--209.
\item Schaefer, M.~B. (1954). \emph{Inter-American Tropical Tuna Commission Bulletin}, 1(2), 23--56.
\item Sch\"ar, S., Pohl, E., and Geldermann, J. (2025). \emph{Journal of Multi-Criteria Decision Analysis}, 32, e70013.
\item Sion, M. (1958). \emph{Pacific Journal of Mathematics}, 8(1), 171--176. von Neumann, J. (1928). \emph{Mathematische Annalen}, 100, 295--320.
\end{itemize}

\end{document}
